\section{Week 1: Statistical Learning Foundations for Large Language Models}

\subsection{Learning Objectives}

By the end of this week, you should be able to:
\begin{itemize}
  \item Formally interpret a Large Language Model (LLM) as a conditional probabilistic model
  \item Re-derive the cross-entropy loss used in next-token prediction
  \item Distinguish probability, confidence, and correctness in model outputs
  \item Explain why LLM behavior must be understood statistically rather than symbolically
  \item Identify early human-factor risks caused by misinterpreting probabilistic outputs
\end{itemize}

\subsection{LLM Components Covered}

This week focuses on the statistical substrate underlying all LLM systems.

\textbf{Modeling and Training}
\begin{itemize}
  \item Conditional probability modeling
  \item Maximum likelihood estimation
  \item Cross-entropy loss
  \item Calibration and uncertainty
\end{itemize}

\textbf{Systems Perspective}
\begin{itemize}
  \item Interpretation of model outputs
  \item Implications of probabilistic behavior for downstream system design
\end{itemize}

\subsection{Conceptual Core: What Kind of Object Is an LLM?}

A Large Language Model defines a probability distribution over sequences of discrete symbols (tokens):

\begin{equation}
P_\theta(x_1, x_2, \ldots, x_T)
\end{equation}

Using the chain rule of probability, this joint distribution can be factorized as:

\begin{equation}
P_\theta(x_1, \ldots, x_T) = \prod_{t=1}^{T} P_\theta(x_t \mid x_1, \ldots, x_{t-1})
\end{equation}

Thus, an LLM is trained to model the conditional distribution of the next token given all previous tokens.

Importantly, an LLM does \emph{not} explicitly store facts, rules, or symbols. Instead, it learns to approximate these conditional distributions from large corpora of text. Many characteristic behaviors of LLMs---including hallucinations, overconfidence, and brittle reasoning failures---follow directly from this probabilistic framing.

\subsection{Maximum Likelihood and Next-Token Prediction}

\subsubsection{Training Objective}

Given a dataset of token sequences
\[
\mathcal{D} = \{x^{(i)}_1, \ldots, x^{(i)}_{T_i}\}_{i=1}^N,
\]
training proceeds by maximizing the log-likelihood of the observed data:

\begin{equation}
\mathcal{L}(\theta) = \sum_{i=1}^{N} \sum_{t=1}^{T_i} 
\log P_\theta(x^{(i)}_t \mid x^{(i)}_{<t})
\end{equation}

Equivalently, this corresponds to minimizing the empirical cross-entropy loss:

\begin{equation}
\mathcal{L}_{CE} = - \mathbb{E}_{x \sim \mathcal{D}}
\left[ \log P_\theta(x_t \mid x_{<t}) \right]
\end{equation}

At each timestep, next-token prediction is mathematically equivalent to multinomial logistic regression over the vocabulary.

\subsubsection{Softmax and Logits}

Let $z \in \mathbb{R}^{|V|}$ denote the vector of logits produced by the model for a vocabulary $V$.  
The softmax function converts logits into probabilities:

\begin{equation}
P(x_t = v \mid x_{<t}) =
\frac{\exp(z_v)}{\sum_{v' \in V} \exp(z_{v'})}
\end{equation}

A critical implication is that probabilities are \emph{relative}. Increasing the probability of one token necessarily decreases the probability of others.

\subsection{Information-Theoretic Interpretation}

The cross-entropy between the true data distribution $p$ and the model distribution $q_\theta$ can be decomposed as:

\begin{equation}
H(p, q_\theta) = H(p) + D_{\mathrm{KL}}(p \parallel q_\theta)
\end{equation}

Since $H(p)$ is fixed, minimizing cross-entropy corresponds to minimizing the Kullback--Leibler divergence between the data distribution and the model distribution.

This highlights an important limitation: minimizing loss ensures distributional similarity, not semantic correctness. As a result, fluent but incorrect outputs can score well under the training objective.

\subsection{Uncertainty, Confidence, and Calibration}

\subsubsection{Probability vs.\ Confidence}

It is essential to distinguish:
\begin{itemize}
  \item \textbf{Probability}: the model-assigned likelihood of a token
  \item \textbf{Confidence}: the sharpness (entropy) of the output distribution
  \item \textbf{Correctness}: correspondence with external ground truth
\end{itemize}

LLMs are frequently miscalibrated, exhibiting high confidence on incorrect outputs and low confidence on rare but correct ones.

\subsubsection{Temperature Scaling}

During inference, logits are often rescaled by a temperature parameter $T$:

\begin{equation}
P_T(v) \propto \exp(z_v / T)
\end{equation}

Lower temperatures ($T < 1$) sharpen distributions, while higher temperatures ($T > 1$) increase entropy. Temperature modifies uncertainty but does not introduce new knowledge.

\subsection{Systems Engineering Implications}

From a systems engineering standpoint, an LLM should be treated as a stochastic subsystem with opaque internal state and non-stationary performance.

\begin{center}
\begin{tabular}{|l|l|}
\hline
\textbf{Assumption} & \textbf{Consequence} \\
\hline
Outputs are probabilistic & Variance must be managed \\
Training data is finite & Distribution shift is inevitable \\
No truth guarantee & External verification is required \\
\hline
\end{tabular}
\end{center}

These properties motivate architectural patterns such as verification layers, retrieval augmentation, and human-in-the-loop review.

\subsection{Human Factors and Failure Modes}

A common early failure mode in LLM systems is automation bias:

\begin{quote}
``The model sounds confident, therefore it must be correct.''
\end{quote}

This arises from a mismatch between human epistemic intuitions and probabilistic text generation. Effective mitigations include structured outputs, confidence visualization, and explicit source attribution.

\subsection{Key References}

\begin{itemize}
  \item \citet{bishop2006prml} — \emph{Pattern Recognition and Machine Learning}
  \item \citet{murphy2012mlpp} — \emph{Machine Learning: A Probabilistic Perspective}
  \item \citet{shannon1948} — \emph{A Mathematical Theory of Communication}
\end{itemize}

\subsection{Recommended University Lectures}

\begin{itemize}
  \item Caltech CS156 — \emph{Learning From Data}  
  \url{https://www.youtube.com/playlist?list=PLD63A284B7615313A}

  \item MIT 6.036 — \emph{Introduction to Machine Learning (OCW)}  
  \url{https://ocw.mit.edu/courses/6-036-introduction-to-machine-learning-fall-2020/}

  \item Stanford CS229 — \emph{Machine Learning} (probability and optimization sections)  
  \url{https://cs229.stanford.edu/}
\end{itemize}

\subsection{Practitioner Lab}

\textbf{Objective:} Implement multinomial logistic regression for a toy sequence task using Python and PyTorch.

\textbf{Tasks:}
\begin{enumerate}
  \item Train a simple next-token classifier
  \item Inspect output probability distributions
  \item Measure accuracy, entropy, and calibration error
\end{enumerate}

\textbf{Deliverable:}
\begin{itemize}
  \item Jupyter notebook
  \item One-page reflection on probabilistic intuition failures
\end{itemize}

\subsection{Week 1 Takeaway}

\begin{quote}
An LLM is not a knowledge base or a reasoner; it is a probabilistic sequence model trained on text.
\end{quote}

All subsequent components---Transformers, alignment, retrieval, tools, and agents---build on this foundation.

